\section{谐振电路与电磁热双向闭环}\label{sec:circuit}
\subsection{实际端口电流}
给定电源/负载、补偿网络拓扑、工作频率和温度状态，
由同一 \(\bm Z_{\rm EM}\) 组合电路阻抗，
在基波线性稳态近似下解
\[
\bm{\mathcal Z}_{\rm circuit}(\bm Z_{\rm EM},\omega,\mathcal C,\bm u)
\bm I=\bm V.
\]
补偿电容、电源和负载不应计入耦合器内部空间热源。
LLC/CLLC 开关器件、整流器和谐波需要频域多谐波或时域模型；
首版只能对明确指定的基波工作域作保证。

\subsection{闭环热时间步}
在热时间 \(t_n\)，以热状态 \(\bm a_n\) 计算材料
\(\mathcal M(T_n)\)，求同源单位端口电磁状态
\((\bm Z_n,\{\bm F_{p,n}\})=\mathcal E_\theta(\mathcal S_n,\mathcal M_n)\)，
解电路电流 \(\bm I_n=\mathcal C(\bm Z_n,\bm u_n)\)，
再由场叠加得到局部 \(q_{i,n}\) 和热动力变化：
\begin{equation}
\boxed{
\bm a_n\longrightarrow \mathcal M(T_n)
\longrightarrow(\bm Z_n,\bm F_n)
\longrightarrow\bm I_n
\longrightarrow q_{i,n}
\longrightarrow\bm a_{n+1}.
}\label{eq:feedback}
\end{equation}
导体电阻、海水电导率和磁性材料参数随温度变化，
可导致阻抗和谐振工况改变。
在温度强反馈情况下应在同一时间步内做隐式/固定点迭代，
弱耦合可采用误差受控的半隐式积分。
只有走通闭环才是双向电磁热代理。

\subsection{慢热/快电分解}
电磁载波周期通常远小于热常数，
故可采用载波相量和慢热包络联立；
启动、负载突变及高频谐波超出包络假设时要切换适当电路瞬态模型。
海水流动和位姿历程也应由真实时间输入，不可在网络中仅附加静态位置标签。
